\documentclass[12pt,a4paper]{article}

%% ============================================
%% RAPPORT - MÉTHODE DES MOINDRES CARRÉS
%% ============================================

\usepackage[utf8]{inputenc}
\usepackage[T1]{fontenc}
\usepackage[french]{babel}
\usepackage{amsmath,amssymb,amsfonts}
\usepackage{graphicx}
\usepackage{xcolor}
\usepackage{listings}
\usepackage{geometry}
\usepackage{fancyhdr}
\usepackage{titlesec}
\usepackage{hyperref}
\usepackage{tcolorbox}
\usepackage{float}
\usepackage{framed}
\usepackage{mdframed}
\usepackage{tikz}
\usepackage{pgfplots}
\pgfplotsset{compat=1.18}
\usetikzlibrary{positioning, arrows.meta, calc}

\geometry{margin=2.5cm}

%% --- COULEURS ---
\definecolor{bleu}{HTML}{003366}
\definecolor{bleuclair}{HTML}{e6f0fa}
\definecolor{noir}{HTML}{000000}
\definecolor{gris}{HTML}{f5f5f5}
\definecolor{vert}{HTML}{228B22}
\definecolor{rouge}{HTML}{CC0000}

%% --- STYLE DU CODE PYTHON ---
\lstdefinestyle{python}{
    backgroundcolor=\color{gris},
    commentstyle=\color{vert}\itshape,
    keywordstyle=\color{bleu}\bfseries,
    numberstyle=\tiny\color{noir},
    stringstyle=\color{rouge},
    basicstyle=\ttfamily\footnotesize,
    breaklines=true,
    numbers=left,
    numbersep=8pt,
    frame=single,
    rulecolor=\color{noir},
    showstringspaces=false,
    tabsize=4,
    language=Python
}
\lstset{style=python}

%% --- ENCADRÉS ---
\newmdenv[
    linecolor=bleu,
    linewidth=2pt,
    backgroundcolor=bleuclair,
    innertopmargin=10pt,
    innerbottommargin=10pt,
    innerleftmargin=10pt,
    innerrightmargin=10pt,
]{resultatbox}

\newmdenv[
    linecolor=noir,
    linewidth=1pt,
    backgroundcolor=white,
    innertopmargin=10pt,
    innerbottommargin=10pt,
    innerleftmargin=10pt,
    innerrightmargin=10pt,
]{theorembox}

%% --- EN-TÊTE ET PIED DE PAGE ---
\pagestyle{fancy}
\fancyhf{}
\fancyhead[L]{\textcolor{bleu}{Méthode des Moindres Carrés}}
\fancyhead[R]{\textcolor{bleu}{Daoud Hechaïchi}}
\fancyfoot[C]{\thepage}
\renewcommand{\headrulewidth}{0.4pt}
\renewcommand{\headrule}{\hbox to\headwidth{\color{bleu}\leaders\hrule height \headrulewidth\hfill}}

%% --- STYLE DES SECTIONS ---
\titleformat{\section}{\Large\bfseries\color{bleu}}{\thesection}{1em}{}[\color{bleu}\titlerule]
\titleformat{\subsection}{\large\bfseries\color{noir}}{\thesubsection}{1em}{}

\hypersetup{
    colorlinks=true,
    linkcolor=bleu,
    urlcolor=bleu,
}

%% --- Commandes personnalisées ---
\newcommand{\R}{\mathbb{R}}
\newcommand{\norm}[1]{\left\| #1 \right\|}
\newcommand{\T}{^\top}
\newcommand{\pinv}{^\dagger}

\begin{document}

%% --- PAGE DE TITRE ---
\begin{titlepage}
    \centering
    \vspace*{1cm}

    {\Large\textbf{\textcolor{bleu}{ALGÈBRE LINÉAIRE \& MACHINE LEARNING}}}\\[0.3cm]
    {\large Fondements Mathématiques}\\[1.5cm]

    \noindent\rule{\textwidth}{2pt}\\[0.8cm]
    {\Huge\textbf{Méthode des}}\\[0.3cm]
    {\Huge\textbf{Moindres Carrés}}\\[0.5cm]
    {\Large\textit{Pseudo-inverse de Moore-Penrose $A^\dagger$}}\\[0.8cm]
    \noindent\rule{\textwidth}{2pt}\\[2cm]

    {\large\textbf{Résolution de systèmes surdéterminés}}\\[0.3cm]
    {\large\textbf{Solution optimale au sens de la norme $L^2$}}\\[3cm]

    \vfill

    \noindent\fbox{\parbox{0.6\textwidth}{
        \centering
        \vspace{0.5cm}
        {\Large\textbf{Daoud Hechaïchi}}\\[0.3cm]
        {\large Rapport Technique}\\[0.1cm]
        {\large Mathématiques Appliquées}\\[0.5cm]
    }}\\[1.5cm]

    {\large\today}
\end{titlepage}

%% --- TABLE DES MATIÈRES ---
\tableofcontents
\newpage

%% ============================================
%% CONTENU
%% ============================================

\section{Introduction}

La \textbf{méthode des moindres carrés} est l'une des techniques les plus importantes en mathématiques appliquées. Inventée indépendamment par \textbf{Carl Friedrich Gauss} (1795) et \textbf{Adrien-Marie Legendre} (1805), elle permet de trouver la meilleure approximation d'un système d'équations linéaires qui n'a pas de solution exacte.

\subsection{Motivation}

En pratique, on rencontre souvent des systèmes \textbf{surdéterminés} : plus d'équations que d'inconnues.

\textbf{Exemples concrets :}
\begin{itemize}
    \item Ajuster une droite à un nuage de points expérimentaux
    \item Calibrer un capteur avec de nombreuses mesures redondantes
    \item Reconstruire une image à partir de données incomplètes
    \item Prédire une variable à partir de plusieurs observations
\end{itemize}

\begin{resultatbox}
\begin{center}
\textbf{\textcolor{bleu}{PROBLÈME CENTRAL}}

Étant donné un système $A\mathbf{x} = \mathbf{b}$ sans solution exacte, trouver le vecteur $\hat{\mathbf{x}}$ qui minimise l'erreur quadratique $\norm{A\mathbf{x} - \mathbf{b}}^2$.
\end{center}
\end{resultatbox}

\subsection{Exemple introductif}

On veut ajuster une droite $y = ax + b$ aux trois points : $(1, 2)$, $(2, 3)$, $(3, 5)$.

Cela donne le système :
\begin{align*}
a + b &= 2 \\
2a + b &= 3 \\
3a + b &= 5
\end{align*}

\textbf{Forme matricielle :}
\[
\underbrace{\begin{pmatrix} 1 & 1 \\ 2 & 1 \\ 3 & 1 \end{pmatrix}}_{A}
\underbrace{\begin{pmatrix} a \\ b \end{pmatrix}}_{\mathbf{x}}
= \underbrace{\begin{pmatrix} 2 \\ 3 \\ 5 \end{pmatrix}}_{\mathbf{b}}
\]

Ici $A$ est une matrice $3 \times 2$ : 3 équations pour 2 inconnues. Le système est \textbf{surdéterminé} et n'a généralement pas de solution exacte.

\newpage
\section{Formulation Mathématique}

\subsection{Le problème des moindres carrés}

\begin{theorembox}
\textbf{Problème des Moindres Carrés}

Soit $A \in \R^{m \times n}$ avec $m > n$ (système surdéterminé) et $\mathbf{b} \in \R^m$.

On cherche $\hat{\mathbf{x}} \in \R^n$ qui minimise :
\begin{equation}
\hat{\mathbf{x}} = \arg\min_{\mathbf{x} \in \R^n} \norm{A\mathbf{x} - \mathbf{b}}_2^2
\end{equation}

où $\norm{\cdot}_2$ désigne la norme euclidienne : $\norm{\mathbf{v}}_2 = \sqrt{v_1^2 + v_2^2 + \cdots + v_n^2}$
\end{theorembox}

\subsection{Développement de la fonction coût}

La fonction à minimiser est :
\begin{align}
f(\mathbf{x}) &= \norm{A\mathbf{x} - \mathbf{b}}_2^2 \\[0.3cm]
&= (A\mathbf{x} - \mathbf{b})\T (A\mathbf{x} - \mathbf{b}) \\[0.3cm]
&= \mathbf{x}\T A\T A \mathbf{x} - 2\mathbf{b}\T A \mathbf{x} + \mathbf{b}\T \mathbf{b}
\end{align}

C'est une fonction \textbf{quadratique convexe} en $\mathbf{x}$ : elle possède un unique minimum global.

\subsection{Condition d'optimalité}

Pour trouver le minimum, on annule le gradient :
\begin{equation}
\nabla f(\mathbf{x}) = 2A\T A \mathbf{x} - 2A\T \mathbf{b} = \mathbf{0}
\end{equation}

\begin{resultatbox}
\begin{center}
\textbf{\textcolor{bleu}{ÉQUATIONS NORMALES}}
\begin{equation}
\boxed{A\T A \hat{\mathbf{x}} = A\T \mathbf{b}}
\end{equation}
\end{center}
\end{resultatbox}

\vspace{0.3cm}

\textbf{Remarque :} Le terme "équations normales" vient du fait que le vecteur erreur $\mathbf{b} - A\hat{\mathbf{x}}$ est \textit{normal} (perpendiculaire) à l'espace colonne de $A$.

\newpage
\section{La Pseudo-inverse de Moore-Penrose}

\subsection{Définition}

Quand $A\T A$ est inversible (ce qui est le cas si les colonnes de $A$ sont linéairement indépendantes), on peut résoudre les équations normales explicitement.

\begin{theorembox}
\textbf{Définition : Pseudo-inverse de Moore-Penrose}

La \textbf{pseudo-inverse} (ou inverse généralisée) de $A$ est :
\begin{equation}
\boxed{A\pinv = (A\T A)^{-1} A\T}
\end{equation}

La solution des moindres carrés s'écrit alors :
\begin{equation}
\boxed{\hat{\mathbf{x}} = A\pinv \mathbf{b} = (A\T A)^{-1} A\T \mathbf{b}}
\end{equation}
\end{theorembox}

\subsection{Pourquoi "pseudo-inverse" ?}

Pour une matrice carrée inversible $A$, on a $A^{-1} A = I$. Pour une matrice rectangulaire, $A\pinv$ joue un rôle similaire :

\begin{itemize}
    \item $A\pinv A = I_n$ (matrice identité $n \times n$)
    \item $A A\pinv \neq I_m$ en général (c'est une matrice de projection)
\end{itemize}

\subsection{Propriétés de Moore-Penrose}

\begin{theorembox}
\textbf{Les 4 conditions de Moore-Penrose}

$A\pinv$ est l'unique matrice satisfaisant :
\begin{enumerate}
    \item $A A\pinv A = A$
    \item $A\pinv A A\pinv = A\pinv$
    \item $(A A\pinv)\T = A A\pinv$ \quad (symétrie)
    \item $(A\pinv A)\T = A\pinv A$ \quad (symétrie)
\end{enumerate}
\end{theorembox}

\subsection{Dimensions des matrices}

\begin{table}[H]
\centering
\begin{tabular}{|l|c|l|}
\hline
\textbf{Matrice} & \textbf{Dimensions} & \textbf{Description} \\
\hline
$A$ & $m \times n$ & Matrice du système ($m > n$) \\
$A\T$ & $n \times m$ & Transposée de $A$ \\
$A\T A$ & $n \times n$ & Matrice de Gram (carrée, symétrique) \\
$(A\T A)^{-1}$ & $n \times n$ & Inverse de la matrice de Gram \\
$A\pinv$ & $n \times m$ & Pseudo-inverse \\
\hline
\end{tabular}
\caption{Dimensions des matrices dans le calcul de $A\pinv$}
\end{table}

\newpage
\section{Interprétation Géométrique}

\subsection{Projection orthogonale}

La solution des moindres carrés a une interprétation géométrique élégante.

\begin{theorembox}
\textbf{Théorème de projection}

$\hat{\mathbf{x}}$ est la solution des moindres carrés si et seulement si le vecteur erreur $\mathbf{r} = \mathbf{b} - A\hat{\mathbf{x}}$ est \textbf{orthogonal} à toutes les colonnes de $A$ :
\begin{equation}
A\T (\mathbf{b} - A\hat{\mathbf{x}}) = \mathbf{0}
\end{equation}

Autrement dit : $A\hat{\mathbf{x}}$ est la \textbf{projection orthogonale} de $\mathbf{b}$ sur l'espace colonne de $A$.
\end{theorembox}

\subsection{Illustration}

\begin{figure}[H]
\centering
\begin{tikzpicture}[scale=1.8]
    % Plan (espace colonne de A)
    \fill[bleuclair, opacity=0.6] (-0.3,-0.3) -- (3.2,-0.3) -- (3.5,1.8) -- (0,1.8) -- cycle;

    % Vecteurs colonnes de A
    \draw[-{Stealth}, thick, bleu] (0,0) -- (2.2,0.4) node[right] {$\mathbf{a}_1$};
    \draw[-{Stealth}, thick, bleu] (0,0) -- (0.6,1.2) node[above] {$\mathbf{a}_2$};

    % Vecteur b (hors du plan)
    \draw[-{Stealth}, very thick, rouge] (0,0) -- (1.6,2.6) node[above right] {$\mathbf{b}$};

    % Projection (dans le plan)
    \draw[-{Stealth}, very thick, vert] (0,0) -- (1.9,1.1) node[below right] {$A\hat{\mathbf{x}}$};

    % Résidu (perpendiculaire)
    \draw[-{Stealth}, thick, orange, dashed] (1.9,1.1) -- (1.6,2.6);
    \node[orange] at (2.1, 1.9) {$\mathbf{r}$};

    % Angle droit
    \draw[thick] (1.9,1.1) -- (1.72,1.25) -- (1.57,1.1);

    % Labels
    \node[bleu] at (1.8, 0.15) {$\text{Im}(A)$};

\end{tikzpicture}
\caption{$A\hat{\mathbf{x}}$ est la projection de $\mathbf{b}$ sur l'espace colonne de $A$. Le vecteur $\mathbf{r} = \mathbf{b} - A\hat{\mathbf{x}}$ est perpendiculaire au plan.}
\end{figure}

\subsection{Matrice de projection}

La matrice qui projette sur l'espace colonne de $A$ est :
\begin{equation}
P = A A\pinv = A(A\T A)^{-1} A\T
\end{equation}

\textbf{Propriétés :}
\begin{itemize}
    \item $P^2 = P$ \quad (idempotence : projeter deux fois = projeter une fois)
    \item $P\T = P$ \quad (symétrie)
    \item $P \mathbf{b} = A\hat{\mathbf{x}}$ \quad (le projeté de $\mathbf{b}$)
\end{itemize}

\newpage
\section{Exemple Numérique Détaillé}

Reprenons l'exemple : ajuster $y = ax + b$ aux points $(1, 2)$, $(2, 3)$, $(3, 5)$.

\subsection{Construction des matrices}

\[
A = \begin{pmatrix} 1 & 1 \\ 2 & 1 \\ 3 & 1 \end{pmatrix}, \quad
\mathbf{b} = \begin{pmatrix} 2 \\ 3 \\ 5 \end{pmatrix}
\]

\subsection{Calcul de $A\T A$}

\[
A\T A = \begin{pmatrix} 1 & 2 & 3 \\ 1 & 1 & 1 \end{pmatrix}
\begin{pmatrix} 1 & 1 \\ 2 & 1 \\ 3 & 1 \end{pmatrix}
= \begin{pmatrix} 1+4+9 & 1+2+3 \\ 1+2+3 & 1+1+1 \end{pmatrix}
= \begin{pmatrix} 14 & 6 \\ 6 & 3 \end{pmatrix}
\]

\subsection{Calcul de $A\T \mathbf{b}$}

\[
A\T \mathbf{b} = \begin{pmatrix} 1 & 2 & 3 \\ 1 & 1 & 1 \end{pmatrix}
\begin{pmatrix} 2 \\ 3 \\ 5 \end{pmatrix}
= \begin{pmatrix} 2+6+15 \\ 2+3+5 \end{pmatrix}
= \begin{pmatrix} 23 \\ 10 \end{pmatrix}
\]

\subsection{Inversion de $A\T A$}

Pour une matrice $2 \times 2$ : $\begin{pmatrix} a & b \\ c & d \end{pmatrix}^{-1} = \frac{1}{ad-bc} \begin{pmatrix} d & -b \\ -c & a \end{pmatrix}$

\[
(A\T A)^{-1} = \frac{1}{14 \cdot 3 - 6 \cdot 6} \begin{pmatrix} 3 & -6 \\ -6 & 14 \end{pmatrix}
= \frac{1}{6} \begin{pmatrix} 3 & -6 \\ -6 & 14 \end{pmatrix}
= \begin{pmatrix} 0.5 & -1 \\ -1 & 7/3 \end{pmatrix}
\]

\subsection{Solution finale}

\[
\hat{\mathbf{x}} = (A\T A)^{-1} A\T \mathbf{b}
= \begin{pmatrix} 0.5 & -1 \\ -1 & 7/3 \end{pmatrix}
\begin{pmatrix} 23 \\ 10 \end{pmatrix}
= \begin{pmatrix} 11.5 - 10 \\ -23 + 70/3 \end{pmatrix}
= \begin{pmatrix} 1.5 \\ 1/3 \end{pmatrix}
\]

\begin{resultatbox}
\begin{center}
\textbf{\textcolor{bleu}{RÉSULTAT}}

La droite des moindres carrés est : $\boxed{y = 1.5x + \frac{1}{3} \approx 1.5x + 0.333}$
\end{center}
\end{resultatbox}

\newpage
\section{Lien avec la SVD}

La \textbf{Décomposition en Valeurs Singulières} (SVD) offre une autre façon de calculer $A\pinv$, plus stable numériquement.

\begin{theorembox}
\textbf{SVD et Pseudo-inverse}

Toute matrice $A \in \R^{m \times n}$ admet une décomposition :
\begin{equation}
A = U \Sigma V\T
\end{equation}

où $U$ et $V$ sont orthogonales, et $\Sigma$ contient les valeurs singulières sur sa diagonale.

La pseudo-inverse est alors :
\begin{equation}
A\pinv = V \Sigma\pinv U\T
\end{equation}

où $\Sigma\pinv$ s'obtient en inversant les valeurs singulières non nulles.
\end{theorembox}

\textbf{Avantage :} Cette méthode fonctionne même quand $A\T A$ est singulière ou mal conditionnée. C'est pourquoi les bibliothèques numériques (NumPy, MATLAB) utilisent la SVD en interne.

\section{Application : Régression Linéaire}

La régression linéaire simple est l'application la plus classique des moindres carrés.

\subsection{Problème}

Étant donné $m$ points $(x_i, y_i)$, trouver la droite $y = ax + b$ qui "passe au mieux" par ces points.

\subsection{Formulation matricielle}

\[
\underbrace{\begin{pmatrix} x_1 & 1 \\ x_2 & 1 \\ \vdots & \vdots \\ x_m & 1 \end{pmatrix}}_{A}
\underbrace{\begin{pmatrix} a \\ b \end{pmatrix}}_{\mathbf{x}}
= \underbrace{\begin{pmatrix} y_1 \\ y_2 \\ \vdots \\ y_m \end{pmatrix}}_{\mathbf{y}}
\]

La solution est : $\begin{pmatrix} \hat{a} \\ \hat{b} \end{pmatrix} = (A\T A)^{-1} A\T \mathbf{y}$

\subsection{Régression polynomiale}

Pour ajuster un polynôme de degré $d$ : $y = a_0 + a_1 x + a_2 x^2 + \cdots + a_d x^d$

\[
A = \begin{pmatrix}
1 & x_1 & x_1^2 & \cdots & x_1^d \\
1 & x_2 & x_2^2 & \cdots & x_2^d \\
\vdots & \vdots & \vdots & \ddots & \vdots \\
1 & x_m & x_m^2 & \cdots & x_m^d
\end{pmatrix}
\]

C'est la \textbf{matrice de Vandermonde}.

\newpage
\section{Implémentation Python}

Voici une implémentation \textbf{from scratch}, sans classes, uniquement avec NumPy.

\subsection{Résolution par les équations normales}

\begin{lstlisting}
import numpy as np

def moindres_carres(A, b):
    """
    Resout min ||Ax - b||^2 par les equations normales.
    """
    # Calcul de A^T A
    ATA = np.dot(A.T, A)

    # Calcul de A^T b
    ATb = np.dot(A.T, b)

    # Resolution du systeme ATA * x = ATb
    x_hat = np.linalg.solve(ATA, ATb)

    return x_hat
\end{lstlisting}

\subsection{Calcul de la pseudo-inverse}

\begin{lstlisting}
def pseudo_inverse(A):
    """
    Calcule A^dagger = (A^T A)^{-1} A^T
    """
    ATA = np.dot(A.T, A)
    ATA_inv = np.linalg.inv(ATA)
    A_pinv = np.dot(ATA_inv, A.T)
    return A_pinv
\end{lstlisting}

\subsection{Exemple complet de regression}

\begin{lstlisting}
import numpy as np

# Donnees
x_data = np.array([1, 2, 3, 4, 5])
y_data = np.array([2.1, 3.9, 6.2, 7.8, 10.1])

# Matrice A : [x, 1]
A = np.column_stack([x_data, np.ones(len(x_data))])

# Solution par equations normales
coeffs = moindres_carres(A, y_data)
a, b = coeffs[0], coeffs[1]

print(f"Droite: y = {a:.4f}x + {b:.4f}")
# Affiche: Droite: y = 1.9800x + 0.0600

# Verification avec pseudo-inverse
A_pinv = pseudo_inverse(A)
coeffs_verif = np.dot(A_pinv, y_data)
print(f"Verification: a = {coeffs_verif[0]:.4f}, b = {coeffs_verif[1]:.4f}")
\end{lstlisting}

\textbf{Note :} NumPy fournit \texttt{np.linalg.pinv(A)} pour la pseudo-inverse via SVD (plus stable), et \texttt{np.linalg.lstsq(A, b)} pour résoudre directement le problème.

\newpage
\section{Régularisation de Tikhonov (Ridge)}

Quand $A\T A$ est mal conditionnée (presque singulière), la solution peut être instable. La \textbf{régularisation} stabilise le problème.

\begin{theorembox}
\textbf{Moindres carrés régularisés (Ridge)}

On ajoute un terme de pénalité sur la norme de $\mathbf{x}$ :
\begin{equation}
\min_{\mathbf{x}} \left( \norm{A\mathbf{x} - \mathbf{b}}_2^2 + \lambda \norm{\mathbf{x}}_2^2 \right)
\end{equation}

La solution devient :
\begin{equation}
\boxed{\hat{\mathbf{x}}_\lambda = (A\T A + \lambda I)^{-1} A\T \mathbf{b}}
\end{equation}
\end{theorembox}

\textbf{Effet du paramètre $\lambda$ :}
\begin{itemize}
    \item $\lambda = 0$ : solution classique des moindres carrés
    \item $\lambda$ petit : légère régularisation
    \item $\lambda$ grand : forte régularisation, solution proche de $\mathbf{0}$
\end{itemize}

Le terme $\lambda I$ garantit que $(A\T A + \lambda I)$ est toujours inversible.

\subsection{Implementation}

\begin{lstlisting}
def moindres_carres_ridge(A, b, lambda_reg):
    """
    Moindres carres avec regularisation Ridge.
    """
    n = A.shape[1]
    ATA = np.dot(A.T, A)
    regularisation = lambda_reg * np.eye(n)
    ATb = np.dot(A.T, b)

    x_hat = np.linalg.solve(ATA + regularisation, ATb)
    return x_hat
\end{lstlisting}

\newpage
\section{Résumé des Formules}

\begin{resultatbox}
\begin{center}
\textbf{\textcolor{bleu}{FORMULES ESSENTIELLES}}
\end{center}

\vspace{0.5cm}

\renewcommand{\arraystretch}{1.8}
\begin{tabular}{|p{4.5cm}|p{8cm}|}
\hline
\textbf{Concept} & \textbf{Formule} \\
\hline
Problème & $\displaystyle \min_{\mathbf{x}} \norm{A\mathbf{x} - \mathbf{b}}_2^2$ \\
\hline
Équations normales & $A\T A \hat{\mathbf{x}} = A\T \mathbf{b}$ \\
\hline
Pseudo-inverse & $A\pinv = (A\T A)^{-1} A\T$ \\
\hline
Solution & $\hat{\mathbf{x}} = A\pinv \mathbf{b} = (A\T A)^{-1} A\T \mathbf{b}$ \\
\hline
Matrice de projection & $P = A A\pinv = A(A\T A)^{-1} A\T$ \\
\hline
Via SVD & $A = U\Sigma V\T \Rightarrow A\pinv = V \Sigma\pinv U\T$ \\
\hline
Ridge (régularisé) & $\hat{\mathbf{x}}_\lambda = (A\T A + \lambda I)^{-1} A\T \mathbf{b}$ \\
\hline
\end{tabular}
\end{resultatbox}

\section{Conclusion}

La méthode des moindres carrés et la pseudo-inverse de Moore-Penrose constituent des outils fondamentaux en mathématiques appliquées :

\begin{enumerate}
    \item \textbf{Universalité} : Applicable à tout système linéaire surdéterminé.

    \item \textbf{Optimalité} : Fournit la meilleure approximation au sens de la norme euclidienne.

    \item \textbf{Interprétation géométrique} : La solution correspond à une projection orthogonale sur l'espace colonne.

    \item \textbf{Fondement du Machine Learning} : La régression linéaire, base de nombreux algorithmes, repose entièrement sur cette théorie.

    \item \textbf{Extensions} : Régularisation (Ridge, LASSO), moindres carrés pondérés.
\end{enumerate}

\vspace{0.5cm}

\begin{resultatbox}
\begin{center}
\textbf{\textcolor{bleu}{À RETENIR}}

La pseudo-inverse $A\pinv = (A\T A)^{-1} A\T$ généralise la notion d'inverse aux matrices non carrées. Elle donne la solution unique de norme minimale au problème des moindres carrés.
\end{center}
\end{resultatbox}

\vspace{1.5cm}

\begin{center}
\fbox{\parbox{0.85\textwidth}{
\centering
\vspace{0.5cm}
\textit{``La méthode des moindres carrés est probablement}\\
\textit{la technique d'estimation la plus importante en statistique.''}\\[0.5cm]
--- \textbf{Carl Friedrich Gauss}, 1809
\vspace{0.5cm}
}}
\end{center}

\vspace{1.5cm}

\begin{flushright}
\textit{Daoud Hechaïchi}\\
Rapport Technique - Mathématiques Appliquées
\end{flushright}

\end{document}
